\section{Conclusión}

El proceso de transformación digital ejecutado en Textiles y Confecciones Atlas E.I.R.L., bajo el marco de cofinanciamiento con ProInnóvate, ha demostrado ser un punto de inflexión fundamental para la sostenibilidad y escalabilidad del negocio. A través de la implementación del ERP Odoo (SaaS) y el desarrollo de un ecosistema omnicanal, la empresa logró superar su estado de madurez digital inicial y erradicar su fuerte dependencia de procesos manuales, registros fragmentados e informales.

\subsection{Cierre y Valoración General}

\begin{itemize}
    \item \textbf{Retorno de Inversión (ROI) Tecnológico:} Más allá de los indicadores financieros directos, el mayor impacto percibido tras estos 6 meses de proyecto es la consecución de "paz mental operativa" y eficiencia estratégica en la toma de decisiones. El nuevo sistema ha permitido:
    \begin{itemize}
        \item \textbf{Recuperar hasta el 70\% del tiempo de la gerencia general}, el cual previamente se consumía en tareas operativas de verificación en almacén, y que hoy se redirecciona íntegramente hacia la planificación de compras y expansión comercial.
        \item \textbf{Eliminar de raíz los "quiebres de stock fantasma"} y los errores matemáticos en el mostrador, garantizando un control riguroso del 100\% de las salidas.
        \item \textbf{Convertir la información en un activo estratégico protegido}, mitigando la vulnerabilidad ante la pérdida de datos en dispositivos locales o la excesiva dependencia en la memoria de antiguos colaboradores.
    \end{itemize}
    \item \textbf{Veredicto Final de la Gerencia:} La dirección ejecutiva de la empresa califica esta implementación como un hito decisivo y de máxima prioridad estratégica, otorgándole un puntaje de evaluación técnica de 4.9 sobre 5 en la hoja de ruta inicial. 
\end{itemize}

En síntesis, este proyecto ha permitido que Textiles y Confecciones Atlas E.I.R.L. trascienda definitivamente su modelo de tienda analógica tradicional, convirtiéndose en una comercializadora moderna, tecnológica e interconectada, plenamente preparada para consolidar su liderazgo y ampliar su red de distribución en todo el sur del Perú (Arequipa, Cusco y Puno), amparada en el control absoluto de sus operaciones desde la nube.
